\chapter{El espacio \texorpdfstring{$L^\infty$}{L infinito} y la dualidad de \texorpdfstring{$L^1$}{L1}}\label{v2:cap11}\label{cap:v2-11}
\section{Una cota que ignora conjuntos nulos}\label{lectura:v2-22}
La norma uniforme cambia si se modifica una función en un solo punto. Para trabajar con clases módulo igualdad casi en todas partes se necesita una cota esencial. Las lecturas 22 y 23 de \cite{valencia2} muestran que esta modificación produce un Banach, pero también que la representación de su dual depende de propiedades de la medida que no aparecen para $1<p<\infty$.
\begin{definicion}
Una función medible $f:X\to\F$ es esencialmente acotada si $\mu(\{|f|>M\})=0$ para algún $M<\infty$. Definimos
\[
 \norm f_\infty=\inf\{M\ge0:\mu(\{|f|>M\})=0\}.
\]
El espacio de estas funciones es $\mathcal L^\infty$; su cociente por igualdad casi en todas partes es $L^\infty$.\index{supremo esencial}
\end{definicion}
Siempre $\norm f_\infty\le\sup_X|f|$. La desigualdad puede ser estricta: la función de Thomae en $[0,1]$, igual a $1/q$ en los racionales irreducibles $a/q$ y cero en irracionales, tiene supremo uno y supremo esencial cero para Lebesgue. En cambio, para una función continua en $\R$ ambos supremos coinciden: si $|f(x_0)|>a$, la continuidad mantiene $|f|>a$ en un intervalo abierto de medida positiva.
\begin{proposicion}\label{v2:linfty-cota}
Si $f\in\mathcal L^\infty$, entonces $|f|\le\norm f_\infty$ casi en todas partes. La aplicación $\norm{\cdot}_\infty$ es una seminorma sobre $\mathcal L^\infty$, cuyo núcleo son las funciones nulas casi en todas partes; induce una norma sobre $L^\infty$.
\end{proposicion}
\begin{proof}
Para cada $n$ elegimos una cota esencial $M_n<\norm f_\infty+1/n$. Fuera de la unión numerable de los conjuntos nulos $\{|f|>M_n\}$ se cumple $|f|\le M_n$ para todo $n$, y al pasar al límite $|f|\le\norm f_\infty$. Por ello $|f+g|\le\norm f_\infty+\norm g_\infty$ fuera de un conjunto nulo, lo que da la desigualdad triangular. También $\norm{\alpha f}_\infty\le|\alpha|\norm f_\infty$; para $\alpha\ne0$ se aplica la misma desigualdad a $f=\alpha^{-1}(\alpha f)$ y se obtiene igualdad. El caso $\alpha=0$ se trata directamente. Finalmente, norma cero equivale a $\mu(\{|f|>0\})=0$. Todas estas afirmaciones son invariantes al cambiar el representante en un conjunto medible nulo.
\end{proof}
\begin{teorema}\label{v2:linfty-banach}
$L^\infty$ es un espacio de Banach. Las funciones simples son densas en su norma.
\end{teorema}
\begin{proof}
Sea $(f_n)$ de Cauchy y elijamos representantes. Quitamos el conjunto nulo que es unión de todos los conjuntos
\[
 \{|f_n|>\norm{f_n}_\infty\},\qquad
 \{|f_n-f_m|>\norm{f_n-f_m}_\infty\},\quad n,m\in\N.
\]
En su complemento, la sucesión es uniformemente de Cauchy y uniformemente acotada, porque las normas de una sucesión de Cauchy están acotadas. Tiene allí un límite uniforme $f$, que extendemos por cero al conjunto excepcional. Es medible como límite de las funciones medibles modificadas por cero en ese conjunto. La cota uniforme prueba $f\in\mathcal L^\infty$. Para $n,m$ grandes, $|f_n-f_m|\le\eps$ fuera del conjunto excepcional; al hacer $m\to\infty$, queda $|f_n-f|\le\eps$ allí. Por tanto $\norm{f_n-f}_\infty\le\eps$.

Para densidad, modificamos primero un representante para que $|f|\le M$ en todo $X$. Dividimos el intervalo $[-M,M]$, o el cuadrado $[-M,M]^2$ en el caso complejo, en una cantidad finita de celdas de diámetro menor que $\eps$. Sus preimágenes son medibles; asignamos a cada una un valor de la celda. La función simple resultante $s$ cumple $|f-s|<\eps$ en todo $X$. Es necesario aproximar el valor de $f$, no solamente $|f|$, para conservar los signos y argumentos complejos.
\end{proof}

\section{Ejercicios propuestos: tamaño esencial y límites de exponentes}
\begin{ejercicio}[Lectura 22, ejercicio 1]\label{ej:v2-l22-1}
Demostrar $\norm f_\infty=\inf\{\sup_{x\in A}|f(x)|:A\in\Acal,\ \mu(X\setminus A)=0\}$.
\end{ejercicio}

\begin{ejercicio}[Lectura 22, ejercicio 2]\label{ej:v2-l22-2}
Demostrar que $L^\infty\subset L^1$ si y solo si $\mu(X)<\infty$.
\end{ejercicio}

\begin{ejercicio}[Lectura 22, ejercicio 3]\label{ej:v2-l22-3}
Demostrar que $L^\infty(\R)$ con medida de Lebesgue no es separable, aunque esta medida sea separable.
\end{ejercicio}

\begin{ejercicio}[Lectura 22, ejercicio 4]\label{ej:v2-l22-4}
Si $f\in L^p$, $1\le p\le\infty$, y $g\in L^\infty$, demostrar $fg\in L^p$ y $\norm{fg}_p\le\norm f_p\norm g_\infty$.
\end{ejercicio}

\begin{ejercicio}[Lectura 22, ejercicio 5]\label{ej:v2-l22-5}
Si $\mu(X)<\infty$, demostrar que $f\in L^\infty$ equivale a: (i) $f\in L^p$ para todo $1\le p<\infty$; (ii) $\sup_{p\ge1}\norm f_p<\infty$. En tal caso, demostrar $\lim_{p\to\infty}\norm f_p=\norm f_\infty$.
\end{ejercicio}

\begin{ejercicio}[Lectura 22, ejercicio 6]\label{ej:v2-l22-6}
Si $f\in L^p(\R)$ para todo $p\ge p_0\ge1$ y además $f\in L^\infty(\R)$, ¿se cumple $\norm f_p\to\norm f_\infty$?
\end{ejercicio}

\begin{ejercicio}[Lectura 22, ejercicio 7]\label{ej:v2-l22-7}
Si una sucesión $x$ pertenece a $\ell^p$ para todo $p\ge p_0\ge1$, ¿se cumple $\norm x_p\to\norm x_\infty$?
\end{ejercicio}


\section{Los emparejamientos de los extremos}\label{lectura:v2-23}
Hölder en los extremos es el ejercicio \ref{ej:v2-l22-4} con $p=1$. Define aplicaciones lineales contractivas
\[
 \Phi:L^\infty\longrightarrow(L^1)^*,\quad\Phi(g)(f)=\int fg\,d\mu,
 \qquad
 \Psi:L^1\longrightarrow(L^\infty)^*,\quad\Psi(f)(g)=\int fg\,d\mu.
\]
En general no son ambas isometrías sobreyectivas. Los siguientes ejemplos separan tres dificultades distintas.
\begin{ejemplo}[Pérdida de inyectividad de $\Phi$]
En $\N$, definimos $\mu(A)=\#A$ si $1\notin A$ y $\mu(A)=\infty$ si $1\in A$. Es una medida: un conjunto que contiene $1$ tiene en cualquier partición disjunta exactamente una pieza que lo contiene, y fuera de $1$ se trata de conteo. Una función integrable debe anularse en $1$ y satisfacer $\sum_{n\ge2}|f(n)|<\infty$; por ello $L^1$ es isométrico a $\ell^1$ tras eliminar la primera coordenada. El único conjunto nulo es el vacío, así que $L^\infty=\ell^\infty$. Sin embargo, $g=\ind{\{1\}}$ tiene norma uno y $\Phi(g)=0$.
\end{ejemplo}
\begin{ejemplo}[Pérdida de sobreyectividad de $\Phi$]
En $X=(0,1)$ tomamos la $\sigma$-álgebra numerable-conumerable y medida de conteo. Toda función medible es constante fuera de un conjunto numerable, por el argumento del ejercicio \ref{ej:v2-l21-4}. Una función integrable tiene soporte numerable: para cada $n$, el conjunto donde $|f|\ge1/n$ debe ser finito. Por tanto
\[
 \varphi(f)=\sum_{x\in(0,1)}x f(x)
\]
está bien definida y $|\varphi(f)|\le\norm f_1$. El producto $x f(x)$ es medible porque está soportado en un conjunto numerable. Si $\varphi(f)=\int fg$ para una $g$ medible esencialmente acotada, al probar con $f=\ind{\{x\}}$ obtenemos $g(x)=x$ para todo $x$. Esa función no es medible para esta $\sigma$-álgebra, contradicción.
\end{ejemplo}
\begin{ejemplo}[Pérdida de sobreyectividad de $\Psi$]\label{v2:limite-extendido}
En el subespacio $c\subset\ell^\infty$ de sucesiones convergentes, el funcional $x\mapsto\lim_nx_n$ tiene norma uno. Hahn--Banach lo extiende a un funcional $L\in(\ell^\infty)^*$ de norma uno. Si $L(x)=\sum_nx_ny_n$ para algún $y\in\ell^1$, al evaluar en cada vector canónico $e_n$ se obtiene $y_n=L(e_n)=0$. Pero $L(1,1,\ldots)=1$, contradicción. Así la extensión existe, pero no está representada por una sucesión sumable.
\end{ejemplo}

\section{Semifinitud y localizabilidad}
\begin{definicion}
La medida $\mu$ es \emph{semifinita} si cada conjunto medible de medida infinita contiene uno de medida positiva finita. Un \emph{supremo esencial} de una familia $\mathcal E\subset\Acal$ es un $H\in\Acal$ tal que $\mu(E\setminus H)=0$ para todo $E\in\mathcal E$, y tal que cualquier $G$ con esa propiedad satisface $\mu(H\setminus G)=0$. La medida es \emph{localizable} si es semifinita y toda familia de conjuntos medibles admite un supremo esencial.\index{medida!semifinita}\index{medida!localizable}
\end{definicion}
\begin{proposicion}\label{v2:sigma-localizable}
Toda medida $\sigma$-finita es localizable.
\end{proposicion}
\begin{proof}
Escribamos $X=\bigcup_n E_n$ con $E_n$ disjuntos y de medida finita. Si $\mu(A)>0$, algún $A\cap E_n$ tiene medida positiva; si $\mu(A)=\infty$, esta observación da el subconjunto de medida positiva finita que requiere semifinitud. Definimos una medida finita
\[
 \lambda(B)=\sum_{n\ge1}\frac{2^{-n}}{1+\mu(E_n)}\mu(B\cap E_n).
\]
Tiene exactamente los mismos conjuntos nulos que $\mu$.

Dada $\mathcal E$, sea $\alpha$ el supremo de las medidas $\lambda$ de uniones numerables de miembros de $\mathcal E$, incluyendo la unión vacía. Elegimos uniones $H_n$ que se acerquen a $\alpha$ y ponemos $H=\bigcup_nH_n$. Sigue siendo una unión numerable de miembros de $\mathcal E$ y $\lambda(H)=\alpha$. Para $E\in\mathcal E$, $H\cup E$ es otra unión admisible, luego $\lambda(E\setminus H)=0$. Si $G$ contiene esencialmente a cada miembro, la unión numerable que forma $H$ satisface $\lambda(H\setminus G)=0$. Como los conjuntos nulos coinciden, $H$ es el supremo esencial para $\mu$.
\end{proof}
\begin{proposicion}\label{v2:l1-isometria}
$\Phi:L^\infty\to(L^1)^*$ es una isometría si y solo si $\mu$ es semifinita. Además $\Psi:L^1\to(L^\infty)^*$ siempre es una isometría.
\end{proposicion}
\begin{proof}
Si $\mu$ es semifinita y $0<a<\norm g_\infty$, el conjunto $\{|g|>a\}$ tiene medida positiva y contiene un $B$ de medida positiva finita. La función
$f=\mu(B)^{-1}\ind B\,\overline g/|g|$ tiene norma $L^1$ uno y $\int fg=\mu(B)^{-1}\int_B|g|>a$. Haciendo $a\uparrow\norm g_\infty$ obtenemos la isometría. Si falla semifinitud, hay un $A$ de medida infinita cuyos subconjuntos medibles de medida finita son nulos. Para $f\in L^1$, cada $A\cap\{|f|\ge1/n\}$ es finito en medida y por tanto nulo; así $f=0$ casi en todas partes en $A$. Se sigue $\Phi(\ind A)=0$ aunque $\norm{\ind A}_\infty=1$.

Para $f\in L^1$, la función $g=\overline f/|f|$ donde $f\ne0$, y cero en los demás puntos, tiene norma esencial a lo sumo uno y satisface $\int fg=\int|f|$. Esto da $\norm{\Psi(f)}\ge\norm f_1$; Hölder da la otra desigualdad.
\end{proof}
La afirmación completa es que $\Phi$ es una isometría sobreyectiva exactamente cuando $\mu$ es localizable. La demostración, incluida la necesidad de localizabilidad, se dará en el teorema \ref{v2:dual-l1-general}, una vez construidas las derivadas de Radon--Nikodym. Aquí ya se han probado sus ingredientes relativos a la norma y se han visto por qué no puede omitirse una condición sobre la medida.

\section{Ejercicios propuestos: funcionales y propiedades de la medida}
\begin{ejercicio}[Lectura 23, ejercicio 1]\label{ej:v2-l23-1}
Si $\mu$ es semifinita y $g$ es medible con $fg\in L^1$ para todo $f\in L^1$, demostrar $g\in L^\infty$.
\end{ejercicio}

\begin{ejercicio}[Lectura 23, ejercicio 2]\label{ej:v2-l23-2}
Para $f\in L^1$, construir $g\in L^\infty$ con $\norm g_\infty\le1$ y $\int fg=\norm f_1$.
\end{ejercicio}

\begin{ejercicio}[Lectura 23, ejercicio 3]\label{ej:v2-l23-3}
Si $\mu$ es semifinita y $g\in L^\infty$, demostrar
\[
 \norm g_\infty=\sup_{\norm f_1\le1}\left|\int fg\,d\mu\right|
 =\sup_{\norm f_1\le1}\norm{fg}_1.
\]
En el caso real se puede omitir el valor absoluto en el primer supremo; en el caso complejo un supremo de números complejos sin orden no está definido.
\end{ejercicio}

\begin{ejercicio}[Lectura 23, ejercicio 4]\label{ej:v2-l23-4}
Dar un espacio de medida finita y una $g\in L^\infty$ tal que $\norm{fg}_1<\norm g_\infty$ para toda $f\in L^1$ con $\norm f_1\le1$.
\end{ejercicio}

\begin{ejercicio}[Lectura 23, ejercicio 5]\label{ej:v2-l23-5}
Para un conjunto arbitrario $I$ con todas sus partes medibles y medida de conteo, demostrar $(\ell^1(I))^*\cong\ell^\infty(I)$ isométricamente.
\end{ejercicio}

\begin{ejercicio}[Lectura 23, ejercicio 6]\label{ej:v2-l23-6}
Supóngase $\mu$ $\sigma$-finita y $\varphi\in(L^1)^*$. Para cada $A$ de medida positiva finita supóngase dada $g_A\in L^\infty(A)$ que representa a $\varphi$ sobre funciones soportadas en $A$. Demostrar: (i) $g_A=g_B$ casi en todas partes en $A\cap B$; (ii) existe $g\in L^\infty(X)$ que representa a $\varphi$ en todo $L^1$.
\end{ejercicio}

\begin{ejercicio}[Lectura 23, ejercicio 7]\label{ej:v2-l23-7}
Sean $(X,\Acal,\mu)$, $(Y,\mathcal B,\nu)$ y $\phi:X\to Y$ medible, con $\mu(\phi^{-1}(A))=\nu(A)$ para todo $A\in\mathcal B$. Demostrar que $T_\phi(g)=g\circ\phi$ define una isometría lineal de $L^\infty(Y)$ en $L^\infty(X)$.
\end{ejercicio}

\section{Qué distingue a los exponentes extremos}
$L^\infty$ conserva el control uniforme después de eliminar un único conjunto nulo común a los datos numerables de una prueba. Su dual, sin embargo, contiene funcionales que no son integrales contra funciones de $L^1$. En la dirección contraria, semifinitud permite detectar la norma mediante funciones integrables y localizabilidad permitirá pegar representaciones construidas en conjuntos de medida finita.
